% PAGES: 165-166
% Continues the manual "\textit{Solution:}" block for the T_N tridiagonal
% example (sec05_c2_1.tex, page 164 / folio 148, which ends right after
% equation (5.43), the $a \geq 0$ case). Page 165 (folio 149) opens with the
% $a < 0$ case of the same derivation -- no fresh "\textit{Solution:}" label
% is printed; the closing "\hfill$\square$" appears after (5.44) once the
% two cases are combined. No \begin{...} environment was left open by
% sec05_c2_1.tex.

If $a < 0$, we find from (5.41) and (5.42) that
\[
\lambda_1 \;>\; \lambda_2 \;>\; \dots \;>\; \lambda_N.
\]

From Theorem 5.6, it follows that $T_N$ is a symmetric positive definite matrix if and
only if $\lambda_N > 0$.

Recall that $\cos(\pi - x) = -\cos(x)$ for all $x \in \R$ and therefore
\[
\cos\left(\frac{N\pi}{N+1}\right) \;=\; \cos\left(\pi - \frac{\pi}{N+1}\right) \;=\; -\cos\left(\frac{\pi}{N+1}\right).
\]
Then,
\[
\lambda_N \;=\; d - 2a\cos\left(\frac{N\pi}{N+1}\right) \;=\; d + 2a\cos\left(\frac{\pi}{N+1}\right),
\]
and $\lambda_N > 0$ is equivalent to
\begin{equation}
d \;>\; -2a\cos\left(\frac{\pi}{N+1}\right) \;=\; 2|a|\cos\left(\frac{\pi}{N+1}\right);
\end{equation}
the last equality comes from the fact that $a < 0$ and therefore $|a| = -a$.

The conditions (5.43) and (5.44) for the matrix $T_N$ to be symmetric positive definite
can be written together as
\[
d \;>\; 2|a|\cos\left(\frac{\pi}{N+1}\right),
\]
which is what we wanted to show; see (5.40).
\hfill$\square$
\par\medskip

Although the solutions above are elegant, they use the fact that the eigenvalues of the
matrices $B_N$ and $T_N$ are known in advance. In practice, deciding whether a matrix
is symmetric positive definite is not done by computing matrix eigenvalues
numerically, which could be computationally expensive. Instead, the Cholesky
decomposition algorithm is applied to the matrix: if the algorithm does not break down,
then the matrix is symmetric positive definite, else it is not. The cost of applying
the Cholesky algorithm to an $n \times n$ matrix, and therefore of deciding whether the
matrix is symmetric positive definite, is $\frac{1}{3}n^3 + O(n^2)$; see section 6.1
for details.

Many matrices arising in practice are symmetric and weakly or strictly diagonally
dominant. The results below are thus important in practical applications.

\begin{theorem}
(i) Any strictly diagonally dominant symmetric matrix with positive entries on the main
diagonal is symmetric positive definite.

Equivalently, all the eigenvalues of a strictly diagonally dominant symmetric matrix
are strictly greater than $0$.

(ii) Any weakly diagonally dominant symmetric matrix with positive entries on the main
diagonal is symmetric positive semidefinite.

% Source wording: "are greater than or equal 0" -- missing "to" -- reproduced
% as printed, per the never-correct-the-author rule.
Equivalently, all the eigenvalues of a weakly diagonally dominant symmetric matrix are
greater than or equal $0$.
\end{theorem}

\begin{proof}
Let $\lambda$ be an eigenvalue of $A$. Note that $\lambda$ is a real number, since $A$
is a symmetric matrix; cf. Theorem 5.1.
% Page break folio 149/150 falls here; the \begin{proof} above stays open
% across it within this same file.
From Gershgorin's Theorem, see Theorem 4.6, it follows that there exists $j$, with
$1 \leq j \leq n$, such that
\begin{equation}
|A(j,j) - \lambda| \;\leq\; R_j,
\end{equation}
where $R_j$ is given by (4.28), since $|\lambda - A(j,j)| = |A(j,j) - \lambda|$.

Since $a \leq |a|$ for any $a \in \R$, we find from (5.45) that
\[
A(j,j) - \lambda \;\leq\; |A(j,j) - \lambda| \;\leq\; R_j,
\]
and therefore
\begin{equation}
\lambda \;\geq\; A(j,j) - R_j.
\end{equation}

(i) If $A$ is a strictly diagonally dominant matrix with positive entries on the main
diagonal, then $A(j,j) = |A(j,j)|$, and, from (4.30), it follows that
\begin{equation}
A(j,j) \;>\; R_j.
\end{equation}
From (5.46) and (5.47) we conclude that $\lambda > 0$ for any eigenvalue $\lambda$ of
$A$, and therefore that the matrix $A$ is symmetric positive definite; cf. Theorem 5.6.

(ii) If $A$ is a weakly diagonally dominant symmetric matrix with positive entries on
the main diagonal, then $A(j,j) = |A(j,j)|$, and, from (4.29), it follows that
\begin{equation}
A(j,j) \;\geq\; R_j.
\end{equation}
From (5.46) and (5.48), we conclude that $\lambda \geq 0$ for any eigenvalue $\lambda$
of $A$, and therefore the matrix $A$ is symmetric positive semidefinite; cf. Theorem
5.6.
\end{proof}

\par\medskip\noindent\textit{Example:}\ Show that the matrix
\[
A_1 \;=\; \begin{pmatrix}
3 & -0.2 & 1.25 & -0.35\\
-0.2 & 1.5 & 0.25 & 1\\
1.25 & 0.25 & 2.5 & 0.5\\
-0.35 & 1 & 0.5 & 4
\end{pmatrix}
\]
is symmetric positive definite.

\begin{answerx}
The matrix $A_1$ is symmetric and all its main diagonal entries are positive. It is
also strictly diagonally dominant, since, for every row of $A_1$, the absolute value of
the main diagonal entry from the row is greater than the sum of the absolute values of
all the other entries in that row:
\begin{align*}
3 \;&>\; |-0.2| + 1.25 + |-0.35| \;=\; 1.8;\\
1.5 \;&>\; |-0.2| + 0.25 + 1 \;=\; 1.45;\\
2.5 \;&>\; 1.25 + 0.25 + 0.5 \;=\; 2;\\
4 \;&>\; |-0.35| + 1 + 0.5 \;=\; 1.85.
\end{align*}
Then, from Theorem 5.7, we conclude that the matrix $A_1$ is symmetric positive
definite. \hfill$\square$
\end{answerx}
